\begin{frame}{Introduction}
    The Intermediate Representation (IR) plays a key role in optimizing compilers; its semantics plays a dual role:
    \begin{itemize}
        \item It should be \textbf{tight} enough that source-level programs can be efficiently compiled to IR code
        \item It should be \textbf{weak} enough to allow desirable IR-level optimization, and to facilitate the mapping onto the target instruction set
    \end{itemize}
    \begin{figure}
    \centering
    \scalebox{0.7}{
    \begin{tikzpicture}[
        node distance=0.5cm,
        every node/.style={font=\small},
        box/.style={
            draw,
            rounded corners,
            minimum width=2.6cm,
            minimum height=1cm,
            align=center,
            fill=BerkeleyTableBlue1!6
        },
        irbox/.style={
            draw,
            rounded corners,
            minimum width=2.8cm,
            minimum height=1.1cm,
            align=center,
            fill=UnipiBlue!9,
            very thick
        },
        arrow/.style={-Latex, thick}
    ]
    
    \node[box] (src) {Source code};
    \node[box, right=of src] (front) {Front-end};
    \node[irbox, right=of front] (ir) {Intermediate\\Representation\\\textbf{IR}};
    \node[box, right=of ir] (opt) {Optimizations};
    \node[box, right=of opt] (back) {Codegen};
    
    \draw[arrow] (src) -- (front);
    \draw[arrow] (front) -- (ir);
    \draw[arrow] (ir) -- (opt);
    \draw[arrow] (opt) -- (back);
    
    \end{tikzpicture}}
    \end{figure}
    \pdfpcnote{
Stress the tension: IR semantics must be tight enough for correct compilation
from the source, but weak enough to permit aggressive optimizations. The pipeline
diagram shows where IR sits — it is the interface between the front-end and all
optimization passes. This dual role is the root cause of all the problems
that follow.
}
\end{frame}
